\section{Method}
\label{sec:method}

NetDream is built on one idea: safe autoscaling requires a model of the
future. Instead of reacting to threshold crossings, NetDream learns a
graph world model of a Kubernetes microservice application, uses it to
imagine the consequences of candidate scaling actions, rejects unsafe
futures, and executes the first action of the best safe plan. This
section formalises the control problem
(Section~\ref{sec:formulation}), specifies the graph world model
(Section~\ref{sec:world-model}), describes imagination-based planning
with the safety filter (Section~\ref{sec:planning}), and summarises how
the trained model is deployed on a real Kubernetes cluster
(Section~\ref{sec:deployment}).

\subsection{Problem Formulation}
\label{sec:formulation}

We represent a Kubernetes application as a directed service graph
$\mathcal{G}=(V,E)$, where each node $i\in V$ is a microservice and each
edge $(i,j)\in E$ is a service-call dependency. At control step $t$,
service $i$ has state
\[
x_{t,i} =
[\mathrm{CPU}_{t,i}, \mathrm{memory}_{t,i}, \mathrm{QPS}_{t,i},
\mathrm{latency}_{t,i}, \mathrm{error}_{t,i}, \mathrm{replicas}_{t,i}]
\in \mathbb{R}^{d}.
\]
The full cluster state is
\[
X_t = [x_{t,1},\ldots,x_{t,N}] \in \mathbb{R}^{N\times d}.
\]

The autoscaler chooses a per-service scaling action
\[
A_t=[a_{t,1},\ldots,a_{t,N}], \qquad a_{t,i}\in\{-1,0,+1\},
\]
corresponding to scaling service $i$ down, holding it fixed, or scaling
it up by one replica. After the action is applied, the cluster evolves
to $X_{t+1}$, receives reward $r_t$, and may incur a service-level
violation $c_t\in\{0,1\}$. The true transition process
\[
X_{t+1}\sim T(X_t,A_t;\mathcal{G})
\]
is unknown and difficult to model analytically because it depends on
request routing, queueing, pod startup delays, workload shifts, and
cross-service dependencies.

The objective is to reduce service-level violations while avoiding
unnecessary replica cost. We use a reward that penalizes both resource
usage and violations:
\[
r_t =
-\lambda_{\mathrm{cost}}\mathrm{Cost}(X_t)
-\lambda_{\mathrm{slo}}c_t,
\qquad
\mathrm{Cost}(X_t)
=
\sum_{i=1}^{N}\mathrm{replicas}_{t,i},
\]
where the violation indicator $c_t = \mathds{1}[\,\exists\,i\in V :
\mathrm{latency}_{t+1,i}>\mathrm{SLO}\,]$ flags whether the transition
into $X_{t+1}$ produced an SLO breach at any service.
However, NetDream does not rely on reward optimization alone. It also
learns an explicit constraint predictor and uses it during planning to
discard unsafe futures before they are executed.

\subsection{Graph World Model}
\label{sec:world-model}

NetDream learns an action-conditioned graph world model
\[
f_\theta(\mathcal{G},X_t,A_t)
\rightarrow
(\Delta\hat{X}_t,\hat{r}_t,\hat{c}_t),
\]
where $\Delta\hat{X}_t$ is the predicted state residual, $\hat{r}_t$ is
the predicted reward, and $\hat{c}_t$ is the predicted probability of a
service-level violation. The next state is predicted as
\[
\hat{X}_{t+1}=X_t+\Delta\hat{X}_t .
\]
Residual prediction makes the model focus on short-term changes rather
than reconstructing absolute state, which is important for stable
multi-step rollouts.

For each service, NetDream first encodes the current state and candidate
scaling action. Because $a_{t,i}\in\{-1,0,+1\}$ is categorical, we
one-hot encode it before concatenation:
\[
h^{(0)}_{t,i}
=
\phi_{\mathrm{enc}}\left(\bigl[x_{t,i}\,\bigl\|\,
\mathrm{onehot}(a_{t,i})\bigr]\right),
\qquad
\mathrm{onehot}(a_{t,i})\in\{0,1\}^{3}.
\]
This makes the model counterfactual: it can predict what would happen
under different possible scaling actions, not merely forecast the next
state under the current policy.

NetDream then applies $L$ layers of graph attention over the service-call
topology. The directed call graph $\mathcal{G}$ is symmetrized for
message passing so that downstream load and upstream backpressure can
both flow along an edge; we write $\mathcal{N}(i)$ for the resulting
undirected neighbourhood of service $i$.
\[
\tilde{h}^{(\ell+1)}_{t,i}
=
\mathrm{GAT}^{(\ell)}
\left(
h^{(\ell)}_{t,i},
\{h^{(\ell)}_{t,j}:j\in\mathcal{N}(i)\}
\right),
\qquad
h^{(\ell+1)}_{t,i}
=
h^{(\ell)}_{t,i}+\tilde{h}^{(\ell+1)}_{t,i}.
\]
Message passing lets the model capture how load, latency, and failures
propagate across connected services. Attention is useful because service
dependencies are heterogeneous: scaling one service may strongly affect
some downstream services while barely affecting others.

The final node embeddings are decoded by three heads. The dynamics head
predicts per-service residual state change:
\[
\Delta\hat{x}_{t,i}=\phi_{\mathrm{dyn}}(h^{(L)}_{t,i}).
\]
The reward and constraint heads operate on a graph-level representation
\[
g_t=\mathrm{READOUT}\left(\{h^{(L)}_{t,i}\}_{i=1}^{N}\right),
\]
where $\mathrm{READOUT}$ is mean pooling, yielding
\[
\hat{r}_t=\phi_{\mathrm{rew}}(g_t),
\qquad
\hat{c}_t=\sigma(\phi_{\mathrm{con}}(g_t)).
\]
Thus, the same world model predicts how the system evolves, how costly
the transition is, and whether the future is unsafe.

\paragraph{Training the world model.}
The world model is trained offline from logged Kubernetes transitions
$(\mathcal{G},X_t,A_t,X_{t+1},r_t,c_t)$ by minimising a multi-task
objective:
\[
\mathcal{L}(\theta)
=
\underbrace{
\left\|
f_{\theta}^{\mathrm{dyn}}(\mathcal{G},X_t,A_t)
-
(X_{t+1}-X_t)
\right\|_2^2
}_{\text{state dynamics}}
+
\alpha
\underbrace{
(\hat{r}_t-r_t)^2
}_{\text{reward}}
+
\beta
\underbrace{
\mathrm{BCE}(\hat{c}_t,c_t)
}_{\text{constraint}} .
\]
We use $\alpha=\beta=0.1$, keeping dynamics prediction as the dominant
learning signal while still learning reward and safety estimates for
planning. Instead of learning a policy directly through trial and error,
NetDream first learns how the cluster responds to actions; the learned
model then becomes a simulator for planning. This avoids unsafe online
exploration in the real cluster and makes the controller inspectable:
predictions, rewards, and violation risks are explicit quantities that
can be evaluated before action execution.

\subsection{Imagination-Based Planning}
\label{sec:planning}

At deployment time, NetDream uses the learned world model for
model-predictive control. Given the current state $X_t$, the planner
samples $K$ candidate action sequences of horizon $H$:
\[
\mathbf{A}^{(k)}_{t:t+H-1}
=
(A^{(k)}_t,\ldots,A^{(k)}_{t+H-1}),
\qquad k=1,\ldots,K .
\]
Each sequence is rolled out inside the world model:
\[
\hat{X}^{(k)}_{t+h+1}
=
\hat{X}^{(k)}_{t+h}
+
f_{\theta}^{\mathrm{dyn}}
(\mathcal{G},\hat{X}^{(k)}_{t+h},A^{(k)}_{t+h}),
\]
while also predicting rewards and violation probabilities
$\hat{r}^{(k)}_{t+h}$ and $\hat{c}^{(k)}_{t+h}$. The predicted return of
a candidate plan is
\[
J^{(k)}
=
\sum_{h=0}^{H-1}
\gamma^h \hat{r}^{(k)}_{t+h}.
\]

Before selecting the best plan, NetDream applies an imagination-based
safety filter. A candidate sequence is accepted only if its predicted
violation probability remains below threshold $\tau$ at every imagined
step:
\[
\max_{0\le h < H}\hat{c}^{(k)}_{t+h}\le \tau .
\]
Let $\mathcal{K}_{\mathrm{safe}}$ denote the set of candidates that pass
this filter. NetDream selects
\[
k^*=\arg\max_{k\in\mathcal{K}_{\mathrm{safe}}} J^{(k)}
\]
and executes only the first action $A^{(k^*)}_t$ in the real cluster. If
no candidate is predicted safe, the controller falls back to a
conservative no-op action. After observing the next real cluster state,
NetDream replans. Algorithm~\ref{alg:netdream} summarises the loop.

\begin{algorithm}[t]
\caption{NetDream: Imagination-Based Safe Autoscaling}
\label{alg:netdream}
\begin{algorithmic}[1]
\STATE \textbf{Input:} graph $\mathcal{G}$, state $X_t$, world model
$f_\theta$, horizon $H=5$, candidates $K=64$, threshold $\tau=0.5$,
discount $\gamma=0.99$
\STATE Sample candidate action sequences
$\{\mathbf{A}^{(k)}_{t:t+H-1}\}_{k=1}^{K}$
\FOR{$k=1$ to $K$}
    \STATE Initialize $\hat{X}^{(k)}_t \leftarrow X_t$,
    $J^{(k)}\leftarrow 0$, $s^{(k)}\leftarrow \mathrm{true}$
    \FOR{$h=0$ to $H-1$}
        \STATE Predict
        $(\Delta\hat{X},\hat{r},\hat{c})
        \leftarrow
        f_\theta(\mathcal{G},\hat{X}^{(k)}_{t+h},A^{(k)}_{t+h})$
        \STATE Update imagined state
        $\hat{X}^{(k)}_{t+h+1}\leftarrow
        \hat{X}^{(k)}_{t+h}+\Delta\hat{X}$
        \STATE Update score
        $J^{(k)}\leftarrow J^{(k)}+\gamma^h\hat{r}$
        \IF{$\hat{c}>\tau$}
            \STATE $s^{(k)}\leftarrow \mathrm{false}$
        \ENDIF
    \ENDFOR
\ENDFOR
\STATE $\mathcal{K}_{\mathrm{safe}}\leftarrow
\{k:s^{(k)}=\mathrm{true}\}$
\IF{$\mathcal{K}_{\mathrm{safe}}\neq\emptyset$}
    \STATE $k^*\leftarrow
    \arg\max_{k\in\mathcal{K}_{\mathrm{safe}}}J^{(k)}$
    \STATE Execute $A^{(k^*)}_t$
\ELSE
    \STATE Execute no-op action
\ENDIF
\STATE Observe $X_{t+1}$ and re-plan
\end{algorithmic}
\end{algorithm}

\subsection{Online Deployment on Kubernetes}
\label{sec:deployment}

NetDream is designed to run as a closed-loop controller against a real
Kubernetes cluster. The trained world model and the imagination planner
are wrapped in a small Python service that integrates with the
cluster's standard observability stack; the controller is the same
binary across local development clusters and managed cloud clusters,
differing only in its kubeconfig target.

\paragraph{State estimator.}
At every control step (5\,s in our deployment), the controller queries
Prometheus for per-service metrics (CPU, memory, request rate, p99
latency, error rate) and reads the current replica count from the
Kubernetes API. These signals are assembled into the state matrix
$X_t\in\mathbb{R}^{N\times d}$ used by the world model.

\paragraph{Planner and action application.}
$X_t$ and the static service-call graph $\mathcal{G}$ are passed to
Algorithm~\ref{alg:netdream}, which returns a per-service action vector
$A^{(k^*)}_t \in\{-1,0,+1\}^N$. Each non-zero entry is applied as a
patch to the corresponding \texttt{Deployment}'s \texttt{scale}
sub-resource via the \texttt{apps/v1} API. If
$\mathcal{K}_{\mathrm{safe}}=\emptyset$, the controller emits a no-op
and logs a fallback event.

\paragraph{Closed loop and offline data collection.}
Each (state, action, next state, reward, violation) tuple observed at
runtime is appended to a transition log, which can be used to retrain
or fine-tune the world model offline. This closes the loop between
deployment and learning without requiring any online exploration in
the live cluster. Detailed cluster configuration, controller pseudocode,
and operational metrics are provided in
Appendix~\ref{app:deployment}.

This planning loop is the point where the world model is used. The model
does not merely provide an auxiliary prediction loss; it is queried
online to simulate possible futures, evaluate their cost, estimate their
safety, and choose the next scaling action. NetDream therefore converts
Kubernetes autoscaling from reactive control into world-model-based
planning.
